\scanpage{24}
SYSTÈMES de RACINES (6) \hfill Tribu n° 59

peut pour tout entier $n\geq1$ trouver une suite d'éléments $x_1,\ldots,x_n\in P$ telle que $y_i=x_1+\ldots+x_i\in P$, $(1\leq i\leq n)$. Mais si $n>\operatorname{Card}P$, alors il existe deux indices $i$, $j$, $(1\leq i,j\leq n,\ i\ne j)$, tels que $y_i=y_j$, ce qui contredit (2).

4) Pour démontrer la proposition, il suffit, en vertu du cor. 2 de la Prop. 14, de montrer qu'il existe une base $(e_k)_{1\leq k\leq\ell}$ de $V$, telle que pour la relation d'ordre lexicographique définie par cette base, tout élément de $P$ est $>0$. Procédons par récurrence sur $\ell=\dim V$. Le cas $\ell=1$ résultant de (1), on peut supposer $\ell\geq2$. Soit $z\in V-0$ tel que $(z,x)\geq0$ pour tout $x\in P$, qui existe d'après 3) et soient $W$ l'hyperplan orthogonal à $z$, et $V'$ le sous-espace de $W$ engendré par $R\cap W$. Alors $R\cap W$ est un système de racines dans $V'$ et $P\cap W$ est clos dans $R\cap W$. Par hypothèse de récurrence, il existe une base $(f_i)_{1\leq i\leq\ell'=\dim V'}$ de $V'$ telle que les éléments de $P\cap W$ soient $>0$ pour l'ordre lexicographique défini par cette base. Il est alors clair que la base $(e_i)_{(1\leq i\leq\ell)}$ où $e_1=z$, $e_{i+1}=f_i$, $(i=1,\ldots,e')$, et où $e_j(j>e'+1)$ est dans le supplément orthogonal de $V'$ dans $W$ a les propriétés requises.

\uline{n° 7}, (plus grande racine, ancien n° 6).

Le présent Congrès n'en dit rien.

\uline{n° 8}, (pouahs et copouahs).

Le réseau engendré par $R$ se notera $Q(R)$, et ses éléments s'appelleront les pois radiciels ; idem $Q(R^*)$ ; le réseau associé à $Q(R^*)$ sera noté $P(R)$, et ses éléments seront les poix ; idem $P(R^*)$. Donner les inclusions.

\mspage{23}
